\chapter{Continuous actions: DDPG, TD3 and SAC}
\label{ch:continuous}

In robotics and control the actions are vectors of real numbers (torques,
forces, steering angles). Pacman's five moves are not, so this chapter uses
the classical continuous control benchmarks: swinging up a pendulum,
making a simulated robot walk (the MuJoCo tasks). Policy gradients with
Gaussian policies (\cref{sec:gaussian}) work there, and PPO is a standard
baseline; but they are on-policy and need many samples. The algorithms of
this chapter are \emph{off-policy actor-critic} methods: like DQN they learn
a Q-function from a replay memory, and they replace the maximisation over
actions, impossible on a continuum, by a learnt actor.

\section{Deterministic policy gradients}

Let $\mu_\theta : \States \to \R^k$ be a deterministic policy and $Q^\mu$ its
action-value function. The idea of DDPG is to move the action
$\mu_\theta(s)$ in the direction that increases $Q^\mu(s, \cdot)$.

\begin{theorem}{Deterministic policy gradient (stated without proof)}{dpg}
  Under regularity conditions (Silver et al., 2014), with $J(\theta) =
  \E[\sum_t \gamma^t R_{t+1}]$ for the policy $\mu_\theta$ and $\rho^\mu$ its
  discounted state distribution,
  \begin{equation}
    \label{eq:dpg}
    \nabla_\theta J(\theta) = \E_{s \sim \rho^\mu}\left[ \nabla_\theta
    \mu_\theta(s)\, \nabla_a Q^\mu(s, a) \big|_{a = \mu_\theta(s)}
    \right].
  \end{equation}
\end{theorem}

The formula is the chain rule applied to $Q^\mu(s, \mu_\theta(s))$, with the
remarkable fact (as in the stochastic policy gradient theorem) that the
effect of $\theta$ on the state distribution does not appear. It is also
the limit of the stochastic policy gradient for Gaussian policies whose
variance tends to $0$. Since the gradient involves no expectation over
actions, it can be estimated from states drawn from a replay memory filled
by an exploratory policy: the deterministic policy gradient is naturally
off-policy.

\section{DDPG}

Deep deterministic policy gradient (DDPG, Lillicrap et al., 2016) is
``DQN for continuous actions''. A critic $Q_\phi$ is trained with the
Bellman target of the deterministic policy, using target networks for both
the critic and the actor:
\begin{equation}
  y = r + \gamma (1 - d)\, Q_{\phi^-}\!\left(s', \mu_{\theta^-}(s')\right),
\end{equation}
and the actor follows \eqref{eq:dpg} with the critic in place of $Q^\mu$:
it maximises $\E_s[Q_\phi(s, \mu_\theta(s))]$ by gradient ascent through the
critic. Exploration adds noise to the actions (Gaussian, or an
Ornstein--Uhlenbeck process in the original paper), and the target networks
are updated by Polyak averaging. DDPG works but is brittle: the actor
exploits the errors of the critic, seeking actions where $Q_\phi$ is
overestimated, and the overestimation is propagated by the targets.

\section{TD3}

Twin delayed DDPG (TD3, Fujimoto et al., 2018) adds three corrections.
\begin{enumerate}
  \item \textbf{Clipped double Q-learning.} Two critics $Q_{\phi_1},
    Q_{\phi_2}$ are trained, and the target uses the smaller of their
    target versions:
    \begin{equation}
      y = r + \gamma (1 - d) \min_{i = 1, 2} Q_{\phi_i^-}(s', \tilde a'),
    \end{equation}
    a pessimistic estimate that counters the maximisation bias.
  \item \textbf{Target policy smoothing.} $\tilde a' = \mu_{\theta^-}(s') +
    \mathrm{clip}(\varepsilon, -c, c)$ with $\varepsilon \sim \mathcal{N}(0,
    \sigma^2)$: the target is averaged over nearby actions, so that a narrow
    peak of the critic cannot be exploited.
  \item \textbf{Delayed policy updates.} The actor and the target networks
    are updated once every two critic updates, so that the actor follows a
    critic that has had time to settle.
\end{enumerate}

\begin{proposition}{The minimum of two estimates is not optimistic}{td3-min}
  Let $Q_1 = q + \varepsilon_1$ and $Q_2 = q + \varepsilon_2$ be two
  estimates of a value $q$ with errors of mean zero. Then $\E[\min(Q_1,
  Q_2)] \leq q$, with strict inequality as soon as $\Prob(\varepsilon_1 \neq
  \varepsilon_2) > 0$.
\end{proposition}

\begin{proof}
  $\min(Q_1, Q_2) = \frac{Q_1 + Q_2}{2} - \frac{\abs{Q_1 - Q_2}}{2}$, whose
  expectation is $q - \frac{1}{2} \E\abs{\varepsilon_1 - \varepsilon_2}$.
\end{proof}

The bias is now downwards. An underestimated action is avoided by the
actor and is not explored further, which limits the damage, whereas an
overestimated one is sought and its error propagated: pessimism is the
safer error.

\section{Soft actor-critic}

Soft actor-critic (SAC, Haarnoja et al., 2018) learns a stochastic policy
that maximises the reward \emph{and} its own entropy.

\begin{definition}{Maximum entropy objective}{maxent}
  For a temperature $\alpha > 0$,
  \begin{equation}
    J_\alpha(\pi) = \E_\pi\left[ \sum_t \gamma^t \left( R_{t+1} + \alpha\,
    \mathcal{H}(\pi(\cdot \mid S_t)) \right) \right],
  \end{equation}
  where $\mathcal{H}(\pi(\cdot \mid s)) = -\E_{a \sim \pi(\cdot \mid s)}[\log
  \pi(a \mid s)]$. The \emph{soft} action values of $\pi$ satisfy
  \begin{equation}
    \label{eq:soft-bellman}
    Q^\pi(s, a) = r(s, a) + \gamma\, \E_{s'}\left[ V^\pi(s') \right],
    \qquad V^\pi(s) = \E_{a \sim \pi(\cdot \mid s)}\left[ Q^\pi(s, a) - \alpha
    \log \pi(a \mid s) \right].
  \end{equation}
\end{definition}

The entropy term rewards the policy for keeping several actions possible:
it explores by construction, it learns several ways of solving a task, and
it is more robust to perturbations. We prove the two steps of soft policy
iteration for finitely many actions (the continuous case is analogous with
densities).

\begin{proposition}{Soft policy evaluation}{soft-eval}
  For a fixed policy $\pi$, the operator $(\T^\pi_{\mathrm{soft}} Q)(s, a) =
  r(s, a) + \gamma\, \E_{s'} \E_{a' \sim \pi}[Q(s', a') - \alpha \log \pi(a'
  \mid s')]$ is a $\gamma$-contraction in the sup norm; its fixed point is
  the soft action-value function $Q^\pi$.
\end{proposition}

\begin{proof}
  The entropy terms do not depend on $Q$: they are an additional reward
  $r'(s, a) = -\alpha \gamma\, \E_{s'} \E_{a'}[\log \pi(a' \mid s')]$,
  bounded for finitely many actions. The operator is therefore the Bellman
  expectation operator for action values with reward $r + r'$, a
  $\gamma$-contraction (\cref{lem:contraction}).
\end{proof}

\begin{lemma}{Soft policy improvement}{soft-improvement}
  Let $\pi_{\mathrm{new}}(\cdot \mid s) \propto \exp(Q^{\pi}(s, \cdot) /
  \alpha)$ for every $s$ (the minimiser of $\mathrm{KL}(\pi'(\cdot \mid s)
  \,\|\, \exp(Q^\pi(s, \cdot)/\alpha) / Z(s))$). Then $Q^{\pi_{\mathrm{new}}}
  \geq Q^\pi$.
\end{lemma}

\begin{proof}
  For a fixed state $s$, let $f(\pi') = \E_{a \sim \pi'}[Q^\pi(s, a) - \alpha
  \log \pi'(a \mid s)]$. Writing $\pi^\star(a) = \exp(Q^\pi(s, a)/\alpha) /
  Z(s)$, we have $Q^\pi(s, a) = \alpha \log \pi^\star(a) + \alpha \log Z(s)$,
  hence $f(\pi') = \alpha \log Z(s) - \alpha\, \mathrm{KL}(\pi' \,\|\,
  \pi^\star)$, maximal at $\pi' = \pi^\star = \pi_{\mathrm{new}}(\cdot \mid
  s)$. In particular $f(\pi_{\mathrm{new}}) \geq f(\pi) = V^\pi(s)$. Then, as
  in the proof of \cref{thm:pi-improvement},
  \begin{align*}
    Q^\pi(s, a) &= r(s, a) + \gamma \E_{s'}[V^\pi(s')] \\
    &\leq r(s, a) + \gamma \E_{s'} \E_{a' \sim \pi_{\mathrm{new}}}
      \left[Q^\pi(s', a') - \alpha \log \pi_{\mathrm{new}}(a' \mid s')
      \right] = (\T^{\pi_{\mathrm{new}}}_{\mathrm{soft}} Q^\pi)(s, a),
  \end{align*}
  and iterating the monotone contraction $\T^{\pi_{\mathrm{new}}}_{\mathrm
  {soft}}$ gives $Q^\pi \leq Q^{\pi_{\mathrm{new}}}$.
\end{proof}

As $\alpha \to 0$, $\pi_{\mathrm{new}}$ concentrates on the greedy actions
and soft policy iteration becomes policy iteration. In practice:
\begin{itemize}
  \item two soft critics are trained (clipped double Q, as in TD3) towards
    $y = r + \gamma (1 - d) \left( \min_i Q_{\phi_i^-}(s', \tilde a') -
    \alpha \log \pi_\theta(\tilde a' \mid s') \right)$ with $\tilde a' \sim
    \pi_\theta(\cdot \mid s')$;
  \item the actor minimises $\E_{s} \E_{a \sim \pi_\theta}[\alpha \log
    \pi_\theta(a \mid s) - \min_i Q_{\phi_i}(s, a)]$, the KL divergence of
    the lemma up to a constant, using the \emph{reparametrisation trick}:
    $a = \tanh(\mu_\theta(s) + \sigma_\theta(s) \odot \xi)$ with $\xi \sim
    \mathcal{N}(0, I)$, so that gradients flow through the sampled action;
  \item the temperature is adjusted automatically so that the entropy of
    the policy stays near a target (for example $-k$ for $k$-dimensional
    actions), by gradient descent on $\alpha\, \E[-\log \pi(a \mid s) -
    \bar{\mathcal{H}}]$.
\end{itemize}

\begin{proposition}{Density of a squashed Gaussian}{tanh-density}
  If $u \in \R^k$ has density $\nu$ and $a = \tanh(u)$ componentwise, then
  \begin{equation}
    \log \pi(a) = \log \nu(u) - \sum_{i=1}^{k} \log\left(1 - \tanh^2
    (u_i)\right).
  \end{equation}
\end{proposition}

\begin{proof}
  The map $u \mapsto \tanh(u)$ is a bijection from $\R^k$ to $(-1, 1)^k$
  with diagonal Jacobian of entries $1 - \tanh^2(u_i) > 0$. By the change of
  variables formula, $\pi(a) = \nu(u) / \prod_i (1 - \tanh^2(u_i))$.
\end{proof}

SAC is sample-efficient, stable across hyperparameters, and one of the
standard algorithms for continuous control; a discrete-action variant also
exists, in which the expectations over actions are computed exactly.

\section{Summary}

\begin{itemize}[nosep]
  \item With continuous actions, a learnt actor replaces the maximisation
    over actions; the deterministic policy gradient moves the action along
    the gradient of the critic, and is naturally off-policy.
  \item DDPG is DQN for continuous actions; TD3 makes it reliable with
    clipped double critics (a downward bias, the safer one), target policy
    smoothing and delayed actor updates.
  \item SAC maximises reward plus entropy; soft policy evaluation and
    improvement mirror their classical versions, and with the
    reparametrisation trick and automatic temperature it is a robust
    default for continuous control.
\end{itemize}
